% Part: computability
% Chapter: computability-theory
% Section: prop-reduce

\documentclass[../../../include/open-logic-section]{subfiles}

\begin{document}

\olfileid{cmp}{thy}{ppr}
\olsection{د راکمېدنې خاصيتونه}

که~$A$، $B$ ته راکمېږي، $A \leq_m B$ ليکو۔ دا نښه وړانديز کوي چې
که $A \leq_m B$ وي، نو~$A$ له~$B$ څخه ``زيات ستونزمن نۀ دے''، او
$B$ د~$A$ ``هومره يا ترې زيات ستونزمن دے''؛ او~$\le_m$ د عددونو د
معمولي~$\le$ په شان سټونه (يا د پرېکړې مسئلې) د پېچلتيا له مخې
مرتبوي۔ لاندې دوه قضيې د دې شهود ملاتړ کوي۔ لومړۍ وايي چې~$\le_m$
انتقالي ده۔
\textbf{د ژباړې د نسخې سمون (OLCMP-032):} د سرچينې د دې وروستۍ
جملې په دوو کلمو کښې څرګندې ليکدوديزې تېروتنې وې؛ په طبيعي پښتو کښې
سمه موخه شوې جمله ژباړل شوې ده۔

\begin{prop}
\ollabel{prop:trans-red}
که $A \leq_m B$ او $B \leq_m C$ وي، نو $A \leq_m C$۔
\end{prop}

\begin{proof}
له~$A$ څخه~$B$ ته د راکمونې او له~$B$ څخه~$C$ ته د راکمونې ترکيب
له~$A$ څخه~$C$ ته راکمونه ورکوي۔
\end{proof}

\begin{prob}
\olref{prop:trans-red} په دې ښودلو ثابت کړئ چې که~$f$ او~$g$ په
ترتيب سره له~$A$ څخه~$B$ او له~$B$ څخه~$C$ ته ډېر-پر-يو راکمونې
وي، نو~$\comp{f}{g}$ له~$A$ څخه~$C$ ته ډېر-پر-يو راکمونه ده۔
\end{prob}

\begin{prop}
\ollabel{prop:reduce}
فرض کړئ $A$ او~$B$ هر سټونه دي، او $A \leq_m B$۔
\begin{enumerate}
\item که~$B$ په محاسبوي ډول د شمېر وړ وي، نو~$A$ هم داسې دے۔
\item که~$B$ محاسبه کېدونکے وي، نو~$A$ هم داسې دے۔
\end{enumerate}
\end{prop}

\begin{proof}
$f$ له~$A$ څخه~$B$ ته ډېر-پر-يو راکمونه وي۔ د لومړۍ دعوې دپاره بس
دا وازمويئ چې که~$B$ د جزوي تابع~$g$ د تعريف ساحه وي، نو~$A$ د
~$\comp{f}{g}$ د تعريف ساحه ده:
\begin{align*}
x \in A & \text{ هله او يوازې هله چې } f(x) \in B \\
& \text{ هله او يوازې هله چې }  g(f(x)) \fdefined.
\end{align*}

د دويمې دعوې دپاره راياد کړئ چې که~$B$ محاسبه کېدونکے وي، نو~$B$
او~$\Complement{B}$ په محاسبوي ډول د شمېر وړ دي
(\olref[cmp]{thm:ce-comp})۔ دا ازمويل ستونزمن نۀ دي چې~$f$ له
$\Complement{A}$ څخه~$\Complement{B}$ ته هم ډېر-پر-يو راکمونه ده؛
نو د دې ثبوت د لومړۍ برخې له مخې~$A$ او~$\Complement{A}$ دواړه
!!{computably
enumerable} دي۔ ځکه~$A$ هم د
\olref[cmp]{thm:ce-comp} له مخې محاسبه کېدونکے دے۔ (په بديل ډول،
کتلے شئ چې $\Char{A} = \comp{f}{\Char{B}}$؛ نو که~$\Char{B}$ محاسبه
کېدونکې وي، $\Char{A}$ هم محاسبه کېدونکې ده۔)
\end{proof}

\begin{prob}
فرض کړئ~$f$ له~$A$ څخه~$B$ ته ډېر-پر-يو راکمونه ده۔ وښيئ چې~$f$ له
$\Complement{A}$ څخه~$\Complement{B}$ ته هم ډېر-پر-يو راکمونه ده۔
\end{prob}

\begin{prob}
وښيئ چې که $f\colon \Nat \to \Nat$ ډېر-پر-يو راکمونه وي، نو
$\Char{A} = \comp{f}{\Char{B}}$۔
\textbf{د ژباړې د نسخې سمون (OLCMP-033):} سرچينې دلته د راکمونې
تابع يوازې د لومړي سټ له غړو څخه دويم سټ ته ټايپ کړې وه، حال دا چې
تعريف ئې د ټولو طبيعي عددونو دپاره تابع غواړي او د سټ غړيتوب د
دوه‌اړخيز شرط په وسيله ساتي۔ دلته د تعريف سمه طبيعي-عددونو تعريف
ساحه او هم‌ساحه کارول شوې ده۔
\end{prob}

\begin{digress}
د \emph{ټيورينګ راکمېدنې} په نوم د راکمېدنې يو ډېر عمومي مفهوم په
نورو سياقونو کښې ګټور دے، په ځانګړي ډول د ناپرېکړتيا د پايلو په
ثبوت کښې۔ پام وکړئ چې د \olref[cmp]{cor:comp-k} له مخې د~$K_0$ متمم
~$K_0$ ته نۀ راکمېږي، ځکه په محاسبوي ډول د شمېر وړ نۀ دے۔ خو په
شهودي ډول، که د~$K_0$ په اړه د پوښتنو ځوابونه مو پېژندل، د هغې د
متمم په اړه د پوښتنو ځوابونه به مو هم پېژندل۔ سټ~$A$ هغه وخت~$B$ ته
د ټيورينګ له مخې راکمېدونکے بلل کېږي چې په~$A$ کښې د پوښتنو ځوابونه
د داسې محاسبه کېدونکې کړنلارې په وسيله وټاکل شي چې د~$B$ په اړه
پوښتنې کولے شي۔ دا له ډېر-پر-يو راکمېدنې زيات ازاد دے، ځکه هلته
(1)~يوازې د~$B$ په اړه يوه پوښتنه کولے شئ، او (2)~د ``هو'' ځواب بايد
د~$A$ د پوښتنې په ``هو'' ځواب، او همداسې ``نه'' په ``نه''، بدل شي۔
بيا هم که~$A$، $B$ ته د ټيورينګ له مخې راکمېدونکے وي او~$B$ محاسبه
کېدونکے وي، نو~$A$ هم محاسبه کېدونکے دے (که څه هم، لکه وليدل مو،
ورته خبره د محاسبوي شمېر وړتيا دپاره نۀ پوره کېږي)۔

بايد د راکمېدنې د هغو بېلابېلو مفهومونو په اړه فکر وکړئ چې بحث مو
پرې وکړ، او توپيرونه ئې درک کړئ۔ خو په دې څپرکي کښې به يوازې له
ډېر-پر-يو راکمېدنې سره کار کوو۔ په ضمني ډول، په تېر پراګراف کښې د
راکمېدنې دواړه ډولونه په محاسبوي پېچلتيا کښې ورته بڼې لري، خو دا
زيات شرط ورسره وي چې ټيورينګ ماشينونه په څوحدي وخت کښې وچلېږي: د
ډېر-پر-يو راکمېدنې پېچلتيايي بڼه \emph{کارپ راکمېدنه}، او د ټيورينګ
راکمېدنې پېچلتيايي بڼه \emph{کوک راکمېدنه} بلل کېږي۔
\end{digress}

\end{document}
